\documentclass[11pt]{article}
\usepackage[margin=1in]{geometry}
\usepackage{amsmath, amssymb}
\usepackage{booktabs}
\usepackage{xcolor}

\title{Why the measured bias ratio runs above $r^2$ at large $r$}
\author{}
\date{}

\begin{document}
\maketitle

\section*{The $r^2$ prediction was a leading-order limit}

In \texttt{sigma\_mismatch\_full\_derivation.pdf} we wrote
\begin{equation}
\frac{|\delta h_w|}{|\delta h_z|} \;\simeq\; r^2,
\end{equation}
but this used the approximation $r^2 \ll R^2$ in the curvature denominators. Keeping
\emph{all} terms gives the exact Newton-step ratio.

\section*{The exact Newton-step ratio}

The two Cases differ only in their second-derivative denominators:
\begin{align}
\text{Case A (mismatch $\sigma_w$):}\quad
\mathrm{NLL}''(\nu_t) &= \frac{V_z''}{\sigma_z^2} + r^2\frac{V_w''}{\sigma_w^2}
  \;\simeq\; 2t^2\bigl(R^2 - r^2\bigr),\\[4pt]
\text{Case B (mismatch $\sigma_z$):}\quad
\mathrm{NLL}''(\nu_t) &= r^2\frac{V_z''}{\sigma_z^2} + \frac{V_w''}{\sigma_w^2}
  \;\simeq\; 2t^2\bigl(R^2 r^2 - 1\bigr).
\end{align}
The two gradients (numerators) have equal magnitude $\propto (r^2-1)$, so:
\begin{equation}
\boxed{\;
\frac{|\delta h_w|}{|\delta h_z|} \;=\; \frac{R^2 r^2 - 1}{R^2 - r^2}.
\;}
\label{eq:exact}
\end{equation}

\section*{Expanding the correction}

Factor out $r^2$:
\begin{equation}
\frac{|\delta h_w|}{|\delta h_z|}
\;=\; r^2 \cdot
\underbrace{\frac{1 - 1/(R^2 r^2)}{1 - r^2/R^2}}_{\text{correction factor}}
\;\approx\; r^2 \cdot \frac{1}{\,1 - r^2/R^2\,}.
\label{eq:correction}
\end{equation}
The correction factor is $\geq 1$ and grows monotonically with $r$. For $r\ll R$ it is
$\approx 1$, recovering the $r^2$ scaling. As $r\to R$ it diverges — at $r=R$ the
Case-A curvature vanishes and the bias truly runs away.

\section*{Numerical example with $R_{\rm eff}=20$}

\begin{center}
\begin{tabular}{rrrrr}
\toprule
$r$ & $r^2/R^2$ & correction & predicted ratio & pure $r^2$ \\
\midrule
1.5 & 0.006 & 1.006 &   2.27 &   2.25 \\
2.0 & 0.010 & 1.010 &   4.04 &   4.00 \\
3.0 & 0.022 & 1.023 &   9.21 &   9.00 \\
5.0 & 0.063 & 1.067 &  26.67 &  25.00 \\
10  & 0.250 & \textbf{1.333} & \textbf{133.3} & 100.0 \\
\bottomrule
\end{tabular}
\end{center}

At $r=10$ with $R_{\rm eff}\sim 20$, the measured ratio sits roughly $33\%$ above the
$r^2$ line — exactly the trend you observed.

\section*{Why this happens physically}

In Case A, shrinking $\sigma_{w,\rm fit}$ by a factor $r$ multiplies the
\emph{negative} curvature of the $V_w$-direction by $r^2$. That negative curvature
cancels against the positive $V_z$-curvature. The total curvature is
$\propto(R^2 - r^2)$ — \emph{decreasing} with $r$. As the valley flattens, the
optimum slides further before the (now weaker) restoring force stops it. So
$|\delta h_w|$ grows \emph{faster} than $r^2$.

In Case B, both curvatures contribute with the same sign — the valley only stiffens
as $r$ grows. The bias $|\delta h_z|$ decreases with $r$ at the rate the leading-order
formula predicts. No correction.

\section*{How to extract $R_{\rm eff}$ from the data}

Fit the measured ratios to \eqref{eq:exact} with $R_{\rm eff}$ as a single free
parameter:
\begin{equation}
\text{ratio}(r;\,R_{\rm eff}) = \frac{R_{\rm eff}^2 r^2 - 1}{R_{\rm eff}^2 - r^2}.
\end{equation}
A clean fit with a physically sensible $R_{\rm eff}$ value (consistent with
$\sigma_w/(\nu_t\sigma_z)$ for the simulation) confirms that the deviation from $r^2$
is the predicted curvature-cancellation effect, not anharmonicity or finite-$N$
noise.

\section*{Summary}

\begin{itemize}
\item $r^2$ is the leading-order prediction, valid for $r\ll R$.
\item The exact ratio is $\dfrac{R^2 r^2 - 1}{R^2 - r^2}$, which always lies
\emph{above} $r^2$ and diverges as $r\to R$.
\item Your observed up-curve at $r=10$ is the precursor of the $r=R$ runaway.
\item A one-parameter fit for $R_{\rm eff}$ on the measured ratios will confirm
the mechanism quantitatively.
\end{itemize}

\end{document}
